% =====================================================================
\chapter{半充填でサイトごとの電子数が 1 になる理由}
\label{app:ph}
% =====================================================================

開放端の系では一般に、端から電子の密度の振動(Friedel 振動)が立つ。それでも半充填のハバード模型の基底状態では、端でも中央でも厳密に $\ev{n_i}=1$ になる。その理由は粒子正孔対称性である。

\section{粒子正孔変換}
二部格子(すべての結合が 2 つの副格子 A、B をつなぐ格子)で、副格子の符号 $\varepsilon_i=\pm1$($i\in A$ で $+1$、$i\in B$ で $-1$)を使って
\begin{equation}
  \Pi:\quad c_{i\sigma}\to\varepsilon_i\cd_{i\sigma},\qquad \cd_{i\sigma}\to\varepsilon_ic_{i\sigma}
\end{equation}
と変換する。反交換関係を保つユニタリーな変換である。各項は次のように変わる。
\begin{itemize}
  \item 電子数:$n_{i\sigma}\to1-n_{i\sigma}$。
  \item ホッピング:結合の両端は副格子が違うので $\varepsilon_i\varepsilon_j=-1$。反交換の符号と合わせて $\cd_{i\sigma}c_{j\sigma}\to\cd_{j\sigma}c_{i\sigma}$ となり、運動エネルギーの項は全体として不変。二部格子でないと($t'$ や奇数の環)、ホッピングの符号が変わってしまう。
  \item 相互作用:$Un_{i\up}n_{i\dn}\to U(1-n_{i\up})(1-n_{i\dn})=Un_{i\up}n_{i\dn}-Un_i+U$。
\end{itemize}
まとめると $\Pi^\dagger H\Pi=H-U(\hat N-N_{\rm site})$ で、半充填のセクター $\hat N=N_{\rm site}$ の中では $H$ は不変である。$\Pi$ は $(N_\up,N_\dn)$ を $(N_{\rm site}-N_\up,N_{\rm site}-N_\dn)$ に写すので、半充填 $N_\up=N_\dn=N_{\rm site}/2$ のセクターは自分自身に写る。

\section{基底状態の一意性}
Lieb の定理\cite{lieb1989}によれば、有限で連結な二部格子上のハバード模型($U>0$、電子数が偶数)の基底状態は、スピンの多重項の縮退を除いて一意で、全スピンは $S=\lvert N_A-N_B\rvert/2$ である。$N_A=N_B$ なら完全に一意な一重項になる。本研究の二部格子の系はすべてこの条件を満たす。厳密対角化でも、1 次元開放端($U=4$、$L=4$〜$10$)で基底状態が一意で最初の励起が三重項であること、$U=0$ では縮退が残ることを確かめた。

一意な基底状態は $\Pi$ で自分自身に写り、位相しか付かない:$\Pi\ket{\psi_0}=e^{i\varphi}\ket{\psi_0}$。

\section{証明}
\begin{equation}
  \ev{n_{i\sigma}}=\bra{\psi_0}\Pi^\dagger n_{i\sigma}\Pi\ket{\psi_0}=\bra{\psi_0}(1-n_{i\sigma})\ket{\psi_0}=1-\ev{n_{i\sigma}}
\end{equation}
より $\ev{n_{i\sigma}}=1/2$、$\ev{n_i}=1$ が、すべてのサイトで成り立つ。サイトの位置を使っていないので、境界条件にも $U$ にもよらない。一般に、$\Pi$ で符号を変える演算子(電荷の偏り $n_i-1$、$Z_{i\up}$、$S^z_i$)の期待値は 0 で、符号を変えない演算子(オンサイト $ZZ$、$\Sv_i\cdot\Sv_j$、ホッピング)には制約がない。

\section{開放端の影響の行き先}
半充填での Friedel 振動の波数は $2k_F=\pi$ で、密度の交替 $(-1)^i$ に当たる。これはちょうど上の対称性で禁止される成分である。そのため開放端の影響は密度には現れず、対称性で禁止されない結合の強弱の交替(隣接スピン相関やホッピングの交替)として現れる。

\section{磁化を持つ UHF でも $\ev{n_i}=1$ になる理由}
$\Pi$ はスピンに対して $y$ 軸まわりの $\pi$ 回転として働き、$S^z_i$ の符号を変える。したがって $z$ 方向に磁化した UHF は $\Pi$ の対称性を破る。しかしハバード模型はスピン回転の対称性も持つので、組み合わせ $\tilde\Pi=\Pi R_y(\pi)^{-1}$ も対称性になる。$\tilde\Pi$ ではスピンは変わらず、電荷だけが反転する。磁化した UHF は $\tilde\Pi$ で自分自身に写るので、$\ev{n_i}=1$ を保つ。実際、二部格子の UHF では大きな磁化($\lvert\ev{S^z_i}\rvert$ が 0.44〜0.49)を持ちながら、$\lvert n_i-1\rvert$ は $10^{-15}$ の程度だった。

二部格子でない $4\times3$ のトーラスでは、この対称性がないので UHF の電荷が 6.5\% 偏った。
